\section{Case D Origin indexed provenance and framing}\label{sec:caseD}

\subsection{Receipt objects and a sequential calculation}
This focused case uses O1 receipt content, not the terminal-grade
interpretation of the full applications. Let $T=\List$ and
$\mathcal L_B=\Set/\List(B)$. A grade is $q:Q\to\List(B)$, a native
map $a:Q\to Q'$ satisfies $q'a=q$, and its extension is $d_BQ=B+Q$.
The occurrence map is $\nu=[\operatorname{singleton},q]$. Transport along
$h:B\to C$ retains the actual receipt carrier and replaces $q$ by
$\List(h)q$. Use any admitted O1 context rig, with unit contexts in this
worked calculation.

Take $B=\{a,b\}$, $Q=\{r,s\}$, $q(r)=q(s)=[a,b]$. A second grade over
$B+Q$ has carrier $R=\{u\}$ and
\[
\rho(u)=[\operatorname{inr}r,\operatorname{inl}a,
                         \operatorname{inr}r].
\]
Receipt accumulation retains $Q+R$ and labels
\[
r\mapsto[a,b],\quad s\mapsto[a,b],\quad
u\mapsto[a,b,a,a,b].
\]
The body selecting $u$ in $(B+Q)+R$, followed by the accumulation body
$(B+Q)+R\to B+(Q+R)$, returns the new retained tag $u$. Its arity
shadow is the displayed five-element list. The two occurrences of $r$
are different addressed positions and both expand. The unused tag $s$
remains in the grade. This is a typed two-stage content protocol, not a
claim that its list is a chronological execution schedule.

The noninvertible native endomorphism $a_Q(r)=r$, $a_Q(s)=r$ preserves
$q$. Since it fixes the referenced $r$, the second-stage label is
unchanged under $\List(\id_B+a_Q)$. The accumulated map is
$a_Q+\id_R$: both routes through accumulation map $r,s$ to $r$ and
$u$ to $u$, and retain its five labels. This checks a genuine native
square including the dormant receipt, not only a final-output equality.

For general $\rho:R\to\List(B+Q)$, the new labels are
\[
[q,\ \mu^{\List}\List(\nu)\rho]:Q+R\longrightarrow\List(B).
\]
At a third stage, each base element contributes its singleton list, each
old receipt contributes its declared list, and each new receipt
contributes the concatenation of its referenced contributions.
Associativity of ordered concatenation identifies both bracketings at
every address. Coproduct reassociation keeps each original stage tag,
so this is also equality on grades and their native maps. For a concrete
label change $h:B\to\{c\}$ with $h(a)=h(b)=c$, the accumulated receipt
$u$ has label $[c,c,c,c,c]$; transport and accumulation give that same
label while retaining $r,s,u$ separately.

\subsection{The exact positive framing comparison}
For an ordinary frame $X$, define
\[
S_XQ=X\times Q,\qquad q^X(x,r)=\List(i_x)q(r),
\quad i_x(b)=(x,b).
\]
For $X=\{0,1\}$ the original two-label receipt becomes respectively
$[(0,a),(0,b)]$ and $[(1,a),(1,b)]$, separately for $r$ and $s$.
The framed $u$ has the five entries
$[(x,a),(x,b),(x,a),(x,a),(x,b)]$. The sequential comparison
\[
(X\times Q)+(X\times R)\cong X\times(Q+R)
\]
retains the common frame and original stage tag, using also
$X\times(B+Q)\cong(X\times B)+(X\times Q)$ at the nested endpoint.
It preserves every displayed label, including the unused $s$.
Repeated framing, units and triple accumulation retain the same tuple
of frames and stage tags, proving their coherence directly.

For $k:X\to Y$, the receipt component is $k\times\id_Q$ over
$k\times\id_B$ and obeys
\[
\List(k\times\id_B)q^X(x,r)=q^Y(k(x),r).
\]
This is invertible when $k$ is bijective; for an inclusion $1\to2$ and
nonempty $Q$ it is not surjective. Thus uniform receipt framing gives
strong sequential comparison at fixed frame and coherent bijection
parameter action, but only a lax comparison for arbitrary frame maps.
Even an empty-list receipt retains $x$ in its tag although its arity
contains no label from which $x$ could be recovered.

\subsection{Full Set framing is obstructed}
The O2 exact empty-fibre criterion applies directly. Since
$\List(\varnothing)=\{[]\}$, a grade in $\mathcal L_\varnothing$ is
an arbitrary receipt set labelled entirely by $[]$. Transport along
$\varnothing\to B$ still labels every receipt by $[]$. For nonempty
$B$, the singleton receipt labelled $[b]$ is not isomorphic over
$\List(B)$ to any such grade. Hence this transport is not essentially
surjective and is not an equivalence. O2's criterion therefore excludes
a full pure-compatible Set-product parameter pseudoaction retaining
these grades.

A direct nonempty-frame witness also rules out a different clever recipe.
At $B=1$ take one receipt labelled $[*]$. Invertible parameter naturality
for the two inclusions $j_i:1\to2$ would identify the same globally
framed grade with singleton grades labelled respectively $[0]$ and $[1]$.
No isomorphism in $\Set/\List(2)$ can identify them. This contradiction
uses retained grades and invertible all-map comparisons, not an empty
frame or a defect particular to the uniform recipe. Lax comparisons
avoid the contradiction exactly as above.

\subsection{What the arity shadow loses}
The strong arity shadow sends a body $f:A\to B+Q$ to
$\nu f:A\to\List(B)$ and forgets receipt-grade maps. Bodies returning
$r$ and $s$ become the same list $[a,b]$, and $a_Q$ has the same shadow
as $\id_Q$. Removing the unused $s$ also leaves a body selecting only
$r$ unchanged in the shadow. The sequential list calculation survives,
but native receipt provenance does not; full framing of the shadow is
not full framing of this source.

This proves the focused indexed-grade, native-cell and framing claims.
It supplies no grade-aware O4 interpretation, stochastic O5 certificate,
or recovery theorem. The finite calculation is checked by
\texttt{boundary\_case\_checks.py}; the general framing conclusion is
the stated O2 criterion and the explicit contradiction, not a finite test.

